\begin{tomtat}
	\begin{minipage}[t]{0.3\textwidth}
		\vspace{0pt}
		\begin{tomtat}
			\begin{center}
				\textbf{ĐỀ 11}
			\end{center}
			
		\end{tomtat}
	\end{minipage}
	\hfill
	\begin{minipage}[t]{0.5\textwidth}
		\vspace{0pt}
		\begin{tomtat}
			\begin{center}
				\includegraphics[scale=0.08]{images/11}
			\end{center}
		\end{tomtat}
	\end{minipage}
\end{tomtat}
\textbf{I. PHẦN TRẮC NGHIỆM}
\begin{ex}
	Trong các biểu thức sau đâu là phương trình bậc nhất 1 ẩn:
	\choice{$x^2-2x=1$}{$3x-2=5$}{$\frac{1}{x}-2=3$}{$0x+8=-5$}
	\loigiai{
		Phương trình bậc nhất một ẩn có dạng $ax+b=0$ với $a \ne 0$.
		\begin{itemize}
			\item A. $x^2-2x=1$ là phương trình bậc hai.
			\item B. $3x-2=5 \Leftrightarrow 3x-7=0$. Đây là phương trình bậc nhất một ẩn với $a=3, b=-7$.
			\item C. $\frac{1}{x}-2=3$ không phải là phương trình bậc nhất (có ẩn ở mẫu).
			\item D. $0x+8=-5 \Leftrightarrow 0x+13=0$. Đây không phải là phương trình bậc nhất vì $a=0$. (Phương trình này vô nghiệm).
		\end{itemize}
		Vậy biểu thức $3x-2=5$ là phương trình bậc nhất một ẩn.
		\\ \textbf{ĐÁP ÁN B}
	}
\end{ex}

\begin{ex}
	Căn bậc hai số học của 9 là:
	\choice{$\sqrt{3}$}{3}{-3;3}{-3}
	\loigiai{
		Căn bậc hai số học của một số không âm $a$ là một số không âm $x$ sao cho $x^2=a$.
		Ta cần tìm số $x \ge 0$ sao cho $x^2=9$.
		Ta có $3^2=9$ và $3 \ge 0$.
		Vậy căn bậc hai số học của 9 là 3.
		\\ \textbf{ĐÁP ÁN B}
	}
\end{ex}

\begin{ex}
	Biểu thức $\frac{(a-2b)(a+b)}{a-2b}$ (với $a \ne 2b$) được rút gọn thành:
	\choice{$a-b$}{$a^2-b^2$}{$a+b$}{$\frac{a}{b}$}
	\loigiai{
		Với điều kiện $a \ne 2b$, ta có $a-2b \ne 0$.
		Do đó, ta có thể rút gọn thừa số $(a-2b)$ ở cả tử và mẫu:
		$$ \frac{(a-2b)(a+b)}{a-2b} = a+b $$
		\\ \textbf{ĐÁP ÁN C}
	}
\end{ex}

\begin{ex}
	Đường thẳng nào dưới đây song song với đường thẳng $y=-5x+3$?
	\choice{$y=1-5x$}{$y=5x+2$}{$y=5x+2025$}{$y=-3+5x$}
	\loigiai{
		Hai đường thẳng $y=ax+b$ và $y=a'x+b'$ song song với nhau khi và chỉ khi $a=a'$ và $b \ne b'$.
		Đường thẳng đã cho là $y=-5x+3$, có $a=-5$ và $b=3$.
		Xét các lựa chọn:
		\begin{itemize}
			\item A. $y=1-5x \Leftrightarrow y=-5x+1$. Có $a'=-5, b'=1$. Vì $a=a' (-5=-5)$ và $b \ne b' (3 \ne 1)$, nên đường thẳng này song song.
			\item B. $y=5x+2$. Có $a'=5 \ne a$. Không song song.
			\item C. $y=5x+2025$. Có $a'=5 \ne a$. Không song song.
			\item D. $y=-3+5x \Leftrightarrow y=5x-3$. Có $a'=5 \ne a$. Không song song.
		\end{itemize}
		Vậy đường thẳng $y=1-5x$ song song với $y=-5x+3$.
		\\ \textbf{ĐÁP ÁN A}
	}
\end{ex}

\begin{ex}
	Trong $\triangle ABC$ vuông tại A, $\tan B$ bằng:
	\choice{$\frac{AB}{AC}$}{$\frac{AC}{AB}$}{$\frac{AB}{BC}$}{$\frac{BC}{AC}$}
	\loigiai{
		Trong tam giác $ABC$ vuông tại A:
		Theo định nghĩa tỉ số lượng giác tang, ta có:
		$$ \tan B = \frac{\text{cạnh đối}}{\text{cạnh kề}} = \frac{AC}{AB} $$
		\\ \textbf{ĐÁP ÁN B}
	}
\end{ex}

\begin{ex}
	Một quả bóng World Cup xem như một hình cầu có đường kính là 17cm. Thể tích hình cầu là:
	\choice{$\frac{4913}{6}\pi \text{ cm}^3$}{$289\pi \text{ cm}^3$}{$4913\pi \text{ cm}^3$}{$\frac{4913}{3}\pi \text{ cm}^3$}
	\loigiai{
		Đường kính của hình cầu là $d = 17 \text{ cm}$.
		Bán kính của hình cầu là $R = \frac{d}{2} = \frac{17}{2} \text{ cm}$.
		Thể tích của hình cầu được tính bằng công thức: $V = \frac{4}{3}\pi R^3$.
		\begin{align*}
			V &= \frac{4}{3}\pi \left(\frac{17}{2}\right)^3 \\
			&= \frac{4}{3}\pi \frac{17^3}{2^3} \\
			&= \frac{4}{3}\pi \frac{4913}{8} \\
			&= \frac{4 \cdot 4913}{3 \cdot 8}\pi \\
			&= \frac{4913}{3 \cdot 2}\pi \\
			&= \frac{4913}{6}\pi (\text{cm}^3)
		\end{align*}
		\\ \textbf{ĐÁP ÁN A}
	}
\end{ex}
\begin{ex} % Câu 7
	Gieo một con xúc xắc 50 lần cho kết quả như sau:
	\begin{center}
		\begin{tabular}{|l|c|c|c|c|c|c|}
			\hline
			Số chấm xuất hiện & 1 & 2 & 3 & 4 & 5 & 6 \\
			\hline
			Tần số & 6 & ? & 8 & 8 & 12 & 7 \\
			\hline
		\end{tabular}
	\end{center}
	Tần số xuất hiện của mặt 2 chấm là:
	\choice{7}{8}{9}{10}
	\loigiai{
		Tổng số lần gieo là 50.
		Tổng tần số của các mặt đã biết là:
		Tần số mặt 1 chấm: 6
		Tần số mặt 3 chấm: 8
		Tần số mặt 4 chấm: 8
		Tần số mặt 5 chấm: 12
		Tần số mặt 6 chấm: 7
		Tổng các tần số đã biết = $6 + 8 + 8 + 12 + 7 = 41$.
		Tần số xuất hiện của mặt 2 chấm là:
		$50 - (\text{Tổng các tần số đã biết})$
		$= 50 - 41 = 9$.
		\\ \textbf{ĐÁP ÁN C}
	}
\end{ex}

\begin{ex} % Câu 8
	Một túi đựng 4 viên bi cùng khối lượng và kích thước được đánh số 1; 2; 3; 4. Lấy ngẫu nhiên hai viên bi từ trong túi. Xác suất để tích hai số ghi trên hai viên bi lớn hơn 3 là:
	\choice{$\frac{5}{6}$}{$\frac{1}{6}$}{$\frac{11}{36}$}{$\frac{2}{3}$}
	\loigiai{
		Tập hợp tất cả các kết quả có thể xảy ra (không gian mẫu $\Omega$) khi lấy lần lượt 2 viên bi từ 4 viên (không hoàn lại) là:
		$\Omega = \{$
		(1,2), (1,3), (1,4), 
		(2,1), (2,3), (2,4), 
		(3,1), (3,2), (3,4), 
		(4,1), (4,2), (4,3) 
		$\}$.
		Số phần tử của không gian mẫu là $n(\Omega) = 12$.
		
		Tập hợp các kết quả thuận lợi cho biến cố B: "tích hai số ghi trên hai viên bi lớn hơn 3" là:
		$B = \{$ (1,4), (2,3), (2,4), (3,2), (3,4), (4,1), (4,2), (4,3) $\}$.
		Số kết quả thuận lợi cho biến cố B là $n(B) = 8$.
		
		Xác suất của biến cố B là $P(B) = \frac{n(B)}{n(\Omega)} = \frac{8}{12} = \frac{2}{3}$.
		\\ \textbf{ĐÁP ÁN D}
	}
\end{ex}
\newpage

\textbf{II. PHẦN TỰ LUẬN}
\begin{ex} % Câu 9
	Rút gọn biểu thức $A = \frac{2}{\sqrt{x}-2} : \left(\frac{\sqrt{x}}{x-4} + \frac{1}{\sqrt{x}-2}\right)$ với $x \ge 0; x \ne 4$.
	\loigiai{
		Với $x \ge 0; x \ne 4$, ta có $x-4 = (\sqrt{x}-2)(\sqrt{x}+2)$.
		\begin{align*}
			A &= \frac{2}{\sqrt{x}-2} : \left(\frac{\sqrt{x}}{(\sqrt{x}-2)(\sqrt{x}+2)} + \frac{1(\sqrt{x}+2)}{(\sqrt{x}-2)(\sqrt{x}+2)}\right) \\
			&= \frac{2}{\sqrt{x}-2} : \left(\frac{\sqrt{x}+\sqrt{x}+2}{(\sqrt{x}-2)(\sqrt{x}+2)}\right) \\
			&= \frac{2}{\sqrt{x}-2} : \frac{2\sqrt{x}+2}{(\sqrt{x}-2)(\sqrt{x}+2)} \\
			&= \frac{2}{\sqrt{x}-2} \cdot \frac{(\sqrt{x}-2)(\sqrt{x}+2)}{2\sqrt{x}+2} \\
			&= \frac{2(\sqrt{x}+2)}{2(\sqrt{x}+1)} \\
			&= \frac{\sqrt{x}+2}{\sqrt{x}+1}
		\end{align*}
		Vậy $A = \frac{\sqrt{x}+2}{\sqrt{x}+1}$.
	}
\end{ex}

\begin{ex} % Câu 10 (Hệ phương trình)
	Giải hệ phương trình: $\begin{cases} 3x-2y=5 \\ 2x+3y=12 \end{cases}$
	\loigiai{
		Ta có hệ phương trình:
		$$ \begin{cases} 3x-2y=5 & (1) \\ 2x+3y=12 & (2) \end{cases} $$
		Nhân phương trình (1) với 3 và phương trình (2) với 2:
		$$ \begin{cases} 9x-6y=15 & (3) \\ 4x+6y=24 & (4) \end{cases} $$
		Cộng vế theo vế của phương trình (3) và (4):
		\begin{align*}
			(9x-6y) + (4x+6y) &= 15+24 \\
			13x &= 39 \\
			x &= 3
		\end{align*}
		Thay $x=3$ vào phương trình (1):
		\begin{align*}
			3(3) - 2y &= 5 \\
			9 - 2y &= 5 \\
			-2y &= 5 - 9 \\
			-2y &= -4 \\
			y &= 2
		\end{align*}
		Vậy nghiệm của hệ phương trình là $(x;y) = (3;2)$.
	}
\end{ex}

\begin{ex} % Câu 10 (Phương trình bậc hai)
	Cho phương trình $x^2-8x+12=0$, không giải phương trình hãy tính giá trị biểu thức sau: $A=\frac{1}{x_1}+\frac{1}{x_2}-\frac{-2}{x_1x_2}$.
	\loigiai{
		Phương trình đã cho là $x^2-8x+12=0$.
		Đây là phương trình bậc hai có dạng $ax^2+bx+c=0$ với $a=1, b=-8, c=12$.
		Ta tính biệt thức $\Delta = b^2-4ac = (-8)^2 - 4(1)(12) = 64 - 48 = 16$.
		Vì $\Delta = 16 > 0$ nên phương trình có hai nghiệm phân biệt $x_1, x_2$.
		Vì $c=12 \ne 0$ nên $x_1, x_2 \ne 0$.
		
		Theo hệ thức Vi-ét, ta có:
		$$ \begin{cases} 
			x_1+x_2 = \dfrac{-b}{a} = \dfrac{-(-8)}{1} = 8 \\ 
			x_1x_2 = \dfrac{c}{a} = \dfrac{12}{1} = 12 
		\end{cases} $$
		Biến đổi biểu thức $A$:
		\begin{align*}
			A &= \frac{1}{x_1}+\frac{1}{x_2}-\frac{-2}{x_1x_2} \\
			&= \frac{x_2}{x_1x_2} + \frac{x_1}{x_1x_2} + \frac{2}{x_1x_2} \\
			&= \frac{x_1+x_2+2}{x_1x_2}
		\end{align*}
		Thay giá trị $x_1+x_2 = 8$ và $x_1x_2 = 12$ vào biểu thức $A$:
		\begin{align*}
			A &= \frac{8+2}{12} \\
			&= \frac{10}{12} \\
			&= \frac{5}{6}
		\end{align*}
		Vậy giá trị của biểu thức $A$ là $\frac{5}{6}$.
	}
\end{ex}

\begin{ex} % Câu 11
	Hai tổ sản xuất của một xí nghiệp ngày thứ nhất dệt được 800 mét vải. Ngày thứ hai do cải tiến kỹ thuật tổ I đã dệt được vượt mức 20\%; tổ II đã dệt được vượt mức 15\% so với ngày thứ nhất nên ngày thứ hai cả hai tổ dệt được 945 mét vải. Hỏi ngày thứ nhất mỗi tổ dệt được bao nhiêu mét vải.
	\loigiai{
		Gọi số mét vải tổ I dệt được trong ngày thứ nhất là $x$ (mét).
		Gọi số mét vải tổ II dệt được trong ngày thứ nhất là $y$ (mét).
		Điều kiện: $x > 0, y > 0$.
		
		Ngày thứ nhất cả hai tổ dệt được 800 mét vải, nên ta có phương trình:
		$$ x + y = 800 \quad (1) $$
		
		Ngày thứ hai:
		Tổ I dệt được: $x + 0,20x = 1,20x$ (mét).
		Tổ II dệt được: $y + 0,15y = 1,15y$ (mét).
		Cả hai tổ dệt được 945 mét vải, nên ta có phương trình:
		$$ 1,20x + 1,15y = 945 \quad (2) $$
		
		Từ (1) và (2) ta có hệ phương trình:
		$$ \begin{cases} x + y = 800 \\ 1,20x + 1,15y = 945 \end{cases} $$
		Từ (1) $\Rightarrow y = 800 - x$. Thay vào (2):
		\begin{align*}
			1,20x + 1,15(800-x) &= 945 \\
			1,20x + 920 - 1,15x &= 945 \\
			0,05x &= 945 - 920 \\
			0,05x &= 25 \\
			x &= \frac{25}{0,05} \\
			x &= 500
		\end{align*}
		Thay $x=500$ vào $y = 800-x$:
		$y = 800 - 500 = 300$.
		Cả $x=500$ và $y=300$ đều thỏa mãn điều kiện.
		
		Vậy, ngày thứ nhất tổ I dệt được 500 mét vải và tổ II dệt được 300 mét vải.
	}
\end{ex}

\begin{ex} % Câu 12
	Một nhóm học sinh có 2 học sinh nam (gồm An và Bình) và 3 học sinh nữ (gồm Châu, Dung và Giang). Cô giáo chọn ra ngẫu nhiên hai học sinh. Tính xác suất của biến cố T: "2 học sinh được chọn có cả nam và nữ".
	\loigiai{
		Tổng số học sinh trong nhóm là $2 (\text{nam}) + 3 (\text{nữ}) = 5$ học sinh.
		Ký hiệu các học sinh: An (A), Bình (B), Châu (C), Dung (D), Giang (G).
		
		Tập hợp tất cả các cách chọn ngẫu nhiên 2 học sinh từ 5 học sinh (không gian mẫu $\Omega$). 
		Các cặp có thể chọn là (lấy không thứ tự, vì chọn "hai học sinh"):
		(A,B), (A,C), (A,D), (A,G)
		(B,C), (B,D), (B,G)
		(C,D), (C,G)
		(D,G) \\
		Tập hợp các kết quả thuận lợi cho biến cố T: "2 học sinh được chọn có cả nam và nữ".
		Điều này có nghĩa là chọn 1 học sinh nam và 1 học sinh nữ.
		\begin{itemize}
			\item Chọn An (nam) và 1 nữ (Châu, Dung, hoặc Giang): (A,C), (A,D), (A,G) - 3 cách.
			\item Chọn Bình (nam) và 1 nữ (Châu, Dung, hoặc Giang): (B,C), (B,D), (B,G) - 3 cách.
		\end{itemize}
		Vậy, các kết quả thuận lợi cho biến cố T là: (A,C), (A,D), (A,G), (B,C), (B,D), (B,G).
		Số phần tử của biến cố T là $6$.
		
		Xác suất của biến cố T là:
		$$ P(T) = \frac{n(T)}{n(\Omega)} = \frac{6}{10} = \frac{3}{5} $$
	}
\end{ex}
\begin{ex} % Câu 13
	Cho tam giác $ABC$ nhọn nội tiếp đường tròn $(O)$. Các đường cao $BD$ và $CE$ cắt nhau tại $H$. Gọi $M, N$ lần lượt là trung điểm của $BC$ và $AH$.
	\begin{enumerate}
		\item[1)] (1,0 điểm) Chứng minh tứ giác $BCDE$ nội tiếp.
		\item[2)] (1,0 điểm) Chứng minh $\widehat{NDM} = \widehat{NEM}$. Gọi $K, L$ lần lượt là giao điểm của hai đường thẳng $OM$ và $CE$, $MN$ và $BD$. Chứng minh $KL$ song song với $AC$.
	\end{enumerate}
	\loigiai{
	\begin{center}
		\begin{tikzpicture}[line join = round, line cap=round,>=stealth,font=\footnotesize,scale=1.5] % Increased scale for clarity
			\usetikzlibrary{calc,intersections}
			
			% Khai báo bán kính
			\def\r{3}
			
			% Định nghĩa tọa độ các điểm ban đầu
			\coordinate (O) at (0,0);
			\coordinate (A) at (130:\r);   % Tọa độ cực A(110 độ, r)
			\coordinate (B) at (-160:\r);  % Tọa độ cực B(-160 độ, r)
			\coordinate (C) at (-20:\r);   % Tọa độ cực C(-20 độ, r)
			
			% E là hình chiếu của C lên AB
			\coordinate (E) at ($(A)!(C)!(B)$);
			% D là hình chiếu của B lên AC
			\coordinate (D) at ($(A)!(B)!(C)$);
			
			% H là giao điểm của CE và BD (Trực tâm của tam giác ABC nếu A,B,C không thẳng hàng)
			% Xác định đường thẳng CE và BD để tìm giao điểm H
			\path[name path=lineCE] (C) -- ($(C)!1.5!(E)$); % Kéo dài CE một chút để chắc chắn cắt
			\path[name path=lineBD] (B) -- ($(B)!1.5!(D)$); % Kéo dài BD một chút
			\path[name intersections={of=lineCE and lineBD, by=H}];
			
			% N là trung điểm của AH
			\coordinate (N) at ($(A)!0.5!(H)$);
			% M là trung điểm của BC
			\coordinate (M) at ($(B)!0.5!(C)$);
			% P là trung điểm của ED
			\coordinate (P) at ($(E)!0.5!(D)$);
			
			% L là giao điểm của MN và BD
			\path[name path=lineMN] (M) -- ($(M)!1.5!(N)$); % Kéo dài MN
			\path[name intersections={of=lineMN and lineBD, by=L}];
			
			% K là giao điểm của OM và CE
			\path[name path=lineOM] (O) -- ($(O)!1.5!(M)$); % Kéo dài OM
			\path[name intersections={of=lineOM and lineCE, by=K}];
			
			% Vẽ các đối tượng hình học
			\draw (O) circle (\r);              % Đường tròn (O,r)
			\draw (A) -- (B) -- (C) -- cycle;   % Tam giác ABC
			\draw (B) -- (D);                   % Đoạn thẳng BD (đường cao từ B)
			\draw (C) -- (E);                   % Đoạn thẳng CE (đường cao từ C)
			\draw (A) -- (H);                   % Đoạn thẳng AH (phần đường cao từ A)
			\draw (E) -- (D) -- (N);                   % Đoạn thẳng ED (cạnh của tam giác pedal)
			\draw (M) -- (E);
			\draw (M) -- (D);
			\draw (M) -- (N) -- (E);
			\draw (O) -- (M);                   % Đoạn OM (đường nối tâm với trung điểm BC)
			\draw (K) -- (L);
			
			% Tô màu và đặt nhãn cho các điểm (theo kiểu cũ)
			\fill (O) circle (1pt);
			\fill (A) circle (1pt);
			\fill (B) circle (1pt);
			\fill (C) circle (1pt);
			\fill (D) circle (1pt);
			\fill (E) circle (1pt);
			\fill (H) circle (1pt);
			\fill (M) circle (1pt);
			\fill (N) circle (1pt);
			\fill (P) circle (1pt);
			\fill (K) circle (1pt);
			\fill (L) circle (1pt);
			
			\node[above] at (O) {$O$};
			\node[above] at (A) {$A$};
			\node[below left] at (B) {$B$};
			\node[below right] at (C) {$C$};
			\node[above right] at (D) {$D$};
			\node[above left] at (E) {$E$};
			\node[below] at (H) {$H$};
			\node[below] at (M) {$M$};
			\node[above right] at (N) {$N$};
			\node[below] at (P) {$P$};
			\node[below right] at (K) {$K$};
			\node[right] at (L) {$L$};
			\pic[draw, angle radius=3mm]{right angle = C--E--B};
			\pic[draw, angle radius=3mm]{right angle = C--D--B};
		\end{tikzpicture}
		
	\end{center}
	\noindent
	\textbf{OM cắt CD tại K. Chứng minh $OK \cdot OM = R^2$.} \\
	\\
	Ta có: Tam giác $ODM$ vuông tại $D$ (vì $\widehat{ODM} = 90^\circ$). \\
	Mặt khác: $MC = MD$ (Tính chất hai tiếp tuyến cắt nhau); \\
	$OC = OD = R$. \\
	Suy ra $OM$ là đường trung trực của đoạn thẳng $CD$ suy ra $OM \perp CD$. \\
	\\
	Trong tam giác vuông $ODM$ tại $D$ và tam giác $OKD$ đồng dạng (g-g) nên ta có: \\
	$OD^2 = OK \cdot OM$. \\
	suy ra $OK \cdot OM = R^2$. \\
	\\
	Theo t/c hai tiếp tuyến cắt nhau ta có $MO$ là tia phân giác của góc $PMQ$, mặt \\
	khác $MO \perp PQ$ nên tam giác $PMQ$ cân tại $M$ suy ra $PQ = 2OP$. \\
	\\
	Ta có $S_{PMQ} = \frac{1}{2} MO \cdot PQ = MO \cdot OP$. \\
	\\
	Trong tam giác vuông $OMP$ ta có: $OC \cdot PM = OP \cdot OM$. \\
	$\Rightarrow OC^2 \cdot PM^2 = OP^2 \cdot OM^2$. \\
	\\
	$\frac{1}{OC^2} = \frac{1}{OP^2} + \frac{1}{OM^2}$ suy ra $\frac{1}{R^2} = \frac{1}{OP^2} + \frac{1}{OM^2}$. \\
	\\
	Áp dụng bất đẳng thức Cô si : \\
	$\frac{1}{OP^2} + \frac{1}{OM^2} \ge 2 \sqrt{\frac{1}{OP^2} \cdot \frac{1}{OM^2}} = \frac{2}{OP \cdot OM}$ suy ra $\frac{1}{R^2} \ge \frac{2}{S_{PMQ}}$. \\
	\\
	$S_{PMQ} \ge 2R^2$. Dấu “=” xảy ra khi
	$\begin{cases}
		\frac{1}{OP^2} = \frac{1}{OM^2} \\
		OM \cdot OP = 2R^2
	\end{cases}$
	suy ra $OM = OP = R\sqrt{2}$. \\
	\\
	Vậy $S_{PMQ}$ đạt giá trị nhỏ nhất khi $OM = R\sqrt{2}$. \\
	\\
	% Nội dung mới từ ảnh thứ hai
	Tứ giác $BCDE$ nội tiếp đường tròn đường kính $BC$ nên \\
	$MD = ME = \frac{1}{2} BC$. \\
	\\
	Tứ giác $AFHD$ nội tiếp đường tròn đường kính $AH$ nên \\
	$ND = NE = \frac{1}{2} AH$. \\
	\\
	$\Rightarrow \triangle NDM = \triangle NEM \text{ (c.c.c)}$. \\
	$\Rightarrow \widehat{NDM} = \widehat{NEM}$. \\
	\\
	Vì $ME = MD$ và $NE = ND$ nên $MN$ là đường trung trực của $DE$. \\
	$\Rightarrow MN \perp DE$. \\
	\\
	Gọi $P$ là giao điểm của $MN$ và $DE$. \\
	\\
	Tam giác $LPD$ có $\widehat{PDL} + \widehat{PLD} = 90^\circ \text{ (1)}$. \\
	\\
	Do $OM$ là trung trực của $BC$ nên tam giác $MKC$ vuông tại $M$, có \\
	$\widehat{KCM} + \widehat{CKM} = 90^\circ \text{ (2)}$. \\
	\\
	Tứ giác $BCDE$ nội tiếp nên $\widehat{KCM} = \widehat{PDL} \text{ (3)}$. \\
	\\
	$\text{Từ (1), (2), (3) } \Rightarrow \widehat{PLD} = \widehat{CKM} \text{ mà } \widehat{PLD} = \widehat{MLB} \text{ (đối đỉnh)},$ \\
	$\widehat{CKM} = \widehat{MKB} \text{ (KM là trung trực của BC) } \Rightarrow \widehat{MLB} = \widehat{MKB}$. \\
	\\
	$\Rightarrow \text{Tứ giác } BMKL \text{ nội tiếp.}$ \\
	\\
	$\text{Vì } \widehat{BMK} = 90^\circ \Rightarrow \widehat{BLK} = 90^\circ \Rightarrow KL \perp BD \text{ mà } AC \perp BD \Rightarrow KL \text{ // } AC.$
}

\end{ex}


\begin{ex} % Câu 14
	Cho các số thực dương $a,b,c$ thỏa mãn $a^2+b^2+c^2=1$.
	Tìm giá trị nhỏ nhất của biểu thức:
	$$ M = \frac{\sqrt{ab+2c^2}}{\sqrt{1+ab-c^2}} + \frac{\sqrt{bc+2a^2}}{\sqrt{1+bc-a^2}} + \frac{\sqrt{ca+2b^2}}{\sqrt{1+ca-b^2}} - 2(ab+bc+ca) $$
\end{ex}

\setcounterref{ex}{0}